\begin{center}
\textsc{Gauss} an \textsc{Olbers}.\quad Göttingen, 2.~September 1837.
\end{center}

\par\noindent\hspace*{\parindent}\mbox{.\ .\ .\ .\ .\ .}\quad \textsc{Gerling}, der seine trigonometrischen Messungen in Hessen jetzt be\-endigt hat, hat jetzt noch eine Operation veranstaltet, die zur Bestimmung des Längenunterschiedes zwischen Göttingen und Mannheim dienen soll.
Es werden Signale auf zwei Bergen, Meisner und Feldberg, gegeben; erstere sind in meiner, letztere in der Mannheimer Sternwarte sichtbar; beide aber zugleich auf einem Zwischenberge bei Marburg, wo \textsc{Gerling} mit einem \textsc{Kessel}\-schen Chronometer beobachtet.
Die Signale sind Pulverzeichen bei Nacht und
\par\smallskip\noindent\hfill{\footnotesize 48\textsuperscript{*}}\hspace*{.12\linewidth}\par
\clearpage
\noindent\makebox[0pt][l]{\fontsize{13}{16}\selectfont \textbf{380}}\hfill{\small\textbf{BRIEFWECHSEL MIT GERLING.}}\hfill\null\par\medskip
\noindent heliotropische bei Tage.
Das Wetter ist nicht günstig; hier sind zwar bisher schon ziemlich viele Zeichen von beiderlei Art beobachtet, aber \textsc{Gerling} hatte nach seinem letzten Briefe noch fast nichts vom Feldberge her gesehen.\ \mbox{.\ .\ .\ .\ .\ .}

\begin{center}\rule{.24\linewidth}{.4pt}\end{center}
